\documentclass[../../main.tex]{subfiles}
\begin{document}

{\bf [解答]}

\begin{figure}[htb]
  \centering
  % 手書きスケッチのアングル（俯角 65°，方位角 120°）
  \tdplotsetmaincoords{65}{100}
  \begin{tikzpicture}[tdplot_main_coords, scale=3.8, line join=round, >=stealth, every node/.style={font=\small}]

    % --- 1. 座標定義 ---
    \coordinate (O)      at (0, 0, 0);
    \coordinate (X1)     at (1, 0, 0);
    \coordinate (Y1)     at (0, 1, 0);
    \coordinate (Z1)     at (0, 0, 1);
    \coordinate (Corner) at (0, 1, 1);

    % 特徴点
    \pgfmathsetmacro{\rMid}{1.0/sqrt(2)} % sqrt(2)/2 ≈ 0.7071
    \coordinate (MidArc) at (0, \rMid, \rMid);           % yz面上の y^2+z^2=1 と y=z の交点
    \coordinate (ApexV)  at (\rMid, \rMid, \rMid);        % 3円柱の交点 (sqrt(2)/2, sqrt(2)/2, sqrt(2)/2)

    % ============================================================
    % 1. 円柱 y^2 + z^2 = 1 の側面 (x = 0 から 1 まで) の薄い斜線塗り
    % ============================================================
    \fill[pattern=north east lines, pattern color=cyan!70!black, opacity=0.35]
    plot[domain=0:90, samples=30] (0, {cos(\x)}, {sin(\x)}) -- (1, 0, 1)
    -- plot[domain=90:0, samples=30] (1, {cos(\x)}, {sin(\x)}) -- cycle;

    % ============================================================
    % 2. 奥の壁面 x = 0 における立体 v の断面 (手書きスケッチの斜線部)
    % ============================================================
    \fill[pattern=north east lines, pattern color=blue!85!black, opacity=0.85]
    (MidArc) -- (Corner) -- (Y1)
    -- plot[domain=0:45, samples=25] (0, {cos(\x)}, {sin(\x)})
    -- cycle;

    % ============================================================
    % 3. 立体 v の外側境界面 (0 <= x <= sqrt(2)/2)
    % ============================================================
    \fill[blue!15, opacity=0.5]
    plot[domain=0:\rMid, samples=25] ({\x}, {sqrt(1.0 - \x*\x)}, {\x})
    -- (\rMid, \rMid, \rMid)
    -- plot[domain=\rMid:0, samples=25] ({\x}, {sqrt(1.0 - \x*\x)}, {sqrt(1.0 - \x*\x)})
    -- (0, 1, 1) -- (0, 1, 0) -- cycle;

    % ============================================================
    % 4. 骨格線・稜線の厳密描画
    % ============================================================
    % yz 平面の枠 (x = 0)
    \draw[thick, gray!75!black] (Z1) -- (Corner) -- (Y1);

    % x=1
    \draw[thick, gray!75!black, dashed] (1,0,1) -- (1,0,0) -- (1,1,0);

    % くり抜き円柱 y^2 + z^2 = 1 (x = 0 および x = 1)
    \draw[very thick, blue!90!black, domain=0:90, samples=40]
    plot (0, {cos(\x)}, {sin(\x)});
    \draw[dashed, gray!55, thick, domain=0:90, samples=40]
    plot (1, {cos(\x)}, {sin(\x)});
    \draw[dashed, gray!45] (0, 1, 0) -- (1, 1, 0);
    \draw[dashed, gray!45] (0, 0, 1) -- (1, 0, 1);

    % 底面円弧 x^2 + y^2 = 1 (z = 0)
    \draw[dashed, gray!55, thick, domain=0:90, samples=30]
    plot ({cos(\x)}, {sin(\x)}, 0);

    % 側面円弧 z^2 + x^2 = 1 (y = 0)
    \draw[dashed, gray!55, thick, domain=0:90, samples=30]
    plot ({cos(\x)}, 0, {sin(\x)});

    % (1,0,0) と (0,1,1) を結ぶ屋根稜線 y = z = sqrt(1-x^2)
    \draw[line width=1.8pt, red!85!black, domain=0:1, samples=45]
    plot ({\x}, {sqrt(1.0 - \x*\x)}, {sqrt(1.0 - \x*\x)});

    % 対称面 y = z の破線 (x = 0 の壁面上: MidArc -> Corner)
    \draw[dashed, red!85!black, line width=1.4pt] (MidArc) -- (Corner);

    % 対称面との交線 (緑太線)
    % (0,0,1) -> ApexV
    \draw[line width=1.8pt, teal!85!black, domain=0:\rMid, samples=30]
    plot ({\x}, {\x}, {sqrt(1.0 - \x*\x)});
    % (0,1,0) -> ApexV
    \draw[line width=1.8pt, teal!85!black, domain=0:\rMid, samples=30]
    plot ({\x}, {sqrt(1.0 - \x*\x)}, {\x});

    % ============================================================
    % ★追加★ 赤い点 ApexV を通る 2 つの円柱の四分円
    % ============================================================
    % (A) 円柱 x^2 + y^2 = 1 の水平スライス (z = sqrt(2)/2 = 一定)
    % (cos t, sin t, sqrt(2)/2)  t: 0° -> 90°
    \draw[line width=1.0pt, orange!90!black, dashdotted, domain=0:90, samples=40]
    plot ({cos(\x)}, {sin(\x)}, {\rMid});
    \node[orange!95!black, font=\tiny\bfseries, right]
    at ({cos(20)}, {sin(20)}, {\rMid}) {$x^2+y^2=1$};

    % (B) 円柱 z^2 + x^2 = 1 の鉛直スライス (y = sqrt(2)/2 = 一定)
    % (cos t, sqrt(2)/2, sin t)  t: 0° -> 90°
    \draw[line width=1.0pt, magenta!85!black, dashdotted, domain=0:90, samples=40]
    plot ({cos(\x)}, {\rMid}, {sin(\x)});
    \node[magenta!90!black, font=\tiny\bfseries, left]
    at ({cos(70)}, {\rMid}, {sin(70)}) {$x^2+z^2=1$};

    % 頂点 (x = y = z = sqrt(2)/2)
    \node[circle, fill=red!85!black, inner sep=1.5pt] at (ApexV) {};
    \node[above left=2pt, font=\scriptsize\bfseries, red!90!black] at (ApexV)
    {$\left(\frac{\sqrt{2}}{2}, \frac{\sqrt{2}}{2}, \frac{\sqrt{2}}{2}\right)$};

    % ============================================================
    % 5. 座標軸と目盛り (最前面)
    % ============================================================
    \draw[->, thick, gray!85!black] (0, 0, 0) -- (1.45, 0, 0) node[below left=1pt] {$x$};
    \draw[->, thick, gray!85!black] (0, 0, 0) -- (0, 1.45, 0) node[right=2pt] {$y$};
    \draw[->, thick, gray!85!black] (0, 0, 0) -- (0, 0, 1.45) node[above=2pt] {$z$};
    \node[below left=2pt] at (O) {$O$};

    \node[below left=1pt, font=\scriptsize] at (X1) {$1$};
    \node[below=2pt, font=\scriptsize]      at (Y1) {$1$};
    \node[left=2pt, font=\scriptsize]       at (Z1) {$1$};

    % 引き出し線ラベル
    \draw[->, thick, blue!90!black] (0.1, 1.35, 0.45) to[out=170, in=-15] (0.1, 0.95, 0.52);
    \node[right=2pt, font=\footnotesize\bfseries, blue!90!black] at (0.1, 1.35, 0.45) {立体 $v$ ($z \le y$ 側)};

  \end{tikzpicture}
  \caption{体積を求める領域のうち$x,y,z \ge 0$の部分}
  \label{fig:three_cylinders_apex_slices}
\end{figure}

簡単のため各軸を$1/r$倍して考える．題意の体積を$V$とする．
与えられた領域で各軸を$1/r$倍した
\begin{align}
  \begin{cases}
    &x^2+y^2\le1 \\
    &y^2+z^2\ge1 \\
    &z^2+x^2\le1
  \end{cases} \label{eq:1}
\end{align}
を満たす立体の体積を$V'$とすると，
\begin{align}
  V=r^3V' \label{eq:2}
\end{align}
が成り立つ．
対称性からさらに$x\ge0$，$0\le z\le y$を満たす部分の体積$v$とすると，
\begin{align}
  V'=16v \label{eq:3}
\end{align}
である．\cref{eq:2,eq:3}より
\begin{align}
  V = 16r^3 v \label{eq:4}
\end{align}
である．

$x=k \ge 0$での切断面を考えると，\cref{eq:1}に代入してかつ$0 \le y \le z$より
\begin{align}
  \begin{cases}
    y^2 \le 1-k^2 \\
    z^2 \le 1-k^2 \\
    y^2 + z^2 \ge 1 \\
    0 \le z \le y
  \end{cases}
\end{align}
だから，$(y,z)$の存在条件および$k \ge 0$より
\begin{align}
  \begin{cases}
    0 \le 1-k^2 \\
    1 \le 2-2k^2
  \end{cases}
  \therefore
  0 \le k \le \dfrac{\sqrt{2}}{2}
\end{align}
である．

そこでパラメータ$\theta$を$0\le\theta\le\dfrac{\pi}{4}$として，
$x=k=\sin\theta$での切断面を考える．以下簡単のため$s=\sin\theta, c=\cos\theta$とすると
\cref{eq:1}より
\begin{align}
  \begin{cases}
    y^2\le1-s^2=c^2 \\
    y^2+z^2\ge1 \\
    z^2\le1-s^2=c^2 \\
    0 \le z \le y
  \end{cases}
\end{align}
であって，この領域は下図斜線部である．

% S(\theta)の様子
\begin{figure}[htb]
  \centering
  \begin{tikzpicture}[x=4.6cm, y=4.6cm, >=stealth, every node/.style={font=\small}]

    % --- 1. 定数・厳密パラメータ (theta = 25° の例) ---
    % s = sin(theta) ≈ 0.4226,  c = cos(theta) ≈ 0.9063
    % 点 P: (y, z) = (c, s), 偏角 theta
    % 点 Q: (y, z) = (c, c) (y = z 上の点)
    % 点 R: y = z (偏角 45°) と 単位円 y^2 + z^2 = 1 の交点 => (cos 45°, sin 45°) = (sqrt(2)/2, sqrt(2)/2)
    \def\angT{25}
    \pgfmathsetmacro{\sVal}{sin(\angT)} % ≈ 0.4226
    \pgfmathsetmacro{\cVal}{cos(\angT)} % ≈ 0.9063
    \pgfmathsetmacro{\rCoord}{cos(45)}  % sqrt(2)/2 ≈ 0.7071

    \coordinate (O) at (0, 0);
    \coordinate (P) at (\cVal, \sVal);
    \coordinate (Q) at (\cVal, \cVal);
    \coordinate (R) at (\rCoord, \rCoord);

    % --- 2. 求める領域 S(theta) の斜線ハッチング (z <= y 側) ---
    % 境界: 点 P(c, s) -> 点 Q(c, c) -> 直線 y=z に沿って点 R へ -> 円弧 R から P へ
    % S(theta) = 三角形 OPQ - 扇形 OPR
    \fill[pattern=north east lines, pattern color=blue!75, opacity=0.8]
    (P) -- (Q) -- (R) arc (45:\angT:1.0) -- cycle;

    % --- 3. 座標軸 (手書き図通り: 横軸 y, 縦軸 z) ---
    \draw[->, thick, gray!80!black] (-0.08, 0) -- (1.18, 0) node[right] {$y$};
    \draw[->, thick, gray!80!black] (0, -0.08) -- (0, 1.18) node[above] {$z$};
    \node[below left=2pt] at (O) {$O$};

    % --- 4. 単位円弧 y^2 + z^2 = 1 ---
    \draw[thick, gray!50, domain=0:90, samples=60] plot ({cos(\x)}, {sin(\x)});
    % 扇形の円弧境界 (R から P へ: 赤太線)
    \draw[very thick, red!85!black, domain=\angT:45, samples=30]
    plot ({cos(\x)}, {sin(\x)});

    % --- 5. 境界枠線 (z <= c, y <= c) ---
    \draw[very thick, black] (0, \cVal) -- (\cVal, \cVal) -- (\cVal, 0);
    \draw[dashed, gray!45] (0, 1.0) -- (1.0, 1.0) -- (1.0, 0);

    % --- 6. 動径と直線 ---
    % 直線 y = z (角度 45°: O から Q を結ぶ線)
    \draw[thick, teal!85!black] (O) -- (Q);

    % 動径 OP (角度 theta)
    \draw[very thick, blue!85!black] (O) -- (P);

    % --- 7. 角度弧とラベル ---
    % 角度 theta (OP の偏角)
    \draw[thick, black] (0.28, 0) arc (0:\angT:0.28);
    \node[font=\footnotesize\bfseries] at ({\angT/2}:0.36) {$\theta$};

    % 角度 pi/4 - theta (直線 y=z と OP のなす角)
    \draw[thick, black] (45:0.42) arc (45:\angT:0.42);
    \node[font=\scriptsize\bfseries, left] at (35:0.56) {$\frac{\pi}{4}-\theta$};

    % --- 8. 投影破線と目盛り ---
    % y 軸目盛り
    \draw[dashed, gray!50] (\sVal, 0) -- (\sVal, \cVal);
    \node[below=2pt, font=\scriptsize] at (\sVal, 0) {$s$};
    \node[below=2pt, font=\scriptsize] at (\cVal, 0) {$c$};
    \node[below=2pt, font=\scriptsize] at (1.0, 0) {$1$};

    % z 軸目盛り
    \draw[dashed, gray!50] (0, \sVal) -- (P);
    \node[left=2pt, font=\scriptsize] at (0, \sVal) {$s$};
    \node[left=2pt, font=\scriptsize] at (0, \cVal) {$c$};
    \node[left=2pt, font=\scriptsize] at (0, 1.0) {$1$};

    % 点プロット
    \node[circle, fill, inner sep=1.2pt] at (P) {}; \node[below right=-1pt, font=\scriptsize] at (P) {$P(c, s)$};
    \node[circle, fill, inner sep=1.2pt] at (Q) {}; \node[above right=-1pt, font=\scriptsize] at (Q) {$Q(c, c)$};
    \node[circle, fill=red!85!black, inner sep=1.2pt] at (R) {}; \node[above left=-1pt, font=\scriptsize, red!90!black] at (R) {$R$};

    % 切断面の見出し
    \node[font=\normalsize\bfseries, above right] at (0.45, 0.95) {$x = s$};

  \end{tikzpicture}
  \caption{切断面 $x = s$ における斜線部 $S(\theta)$の様子}
  \label{fig:cross_section_stheta_area_fixed}
\end{figure}

この斜線部の面積$S(\theta)$として，
\begin{align}
  S(\theta)
  &= \triangle OPQ - \text{扇形}OPR
  = \dfrac{1}{2}(c-s)c-\dfrac{1}{2}\left(\dfrac{\pi}{4}-\theta\right)
\end{align}
であるから，これを積分した体積は
\begin{align}
  v
  &=\int_0^{\sqrt{2}/2}S(\theta)ds \\
  &=\int_0^{\pi/4}S(\theta)\dfrac{ds}{d\theta}d\theta \\
  &=\dfrac{1}{2}\int_0^{\pi/4}\left\{(c-s)c^2-c\left(\dfrac{\pi}{4}-\theta\right)\right\}d\theta \label{eq:5}
\end{align}
である．

各項を個別に計算して，
\begin{align}
  (\text{第一項})
  &= \int_0^{\pi/4}c^3d\theta=\left[s-\dfrac{1}{3}s^3\right]_0^{\pi/4}
  =\dfrac{5\sqrt{2}}{12} \\
  %
  (\text{第二項})
  &=\int_0^{\pi/4}sc^2d\theta=\dfrac{-1}{3}[c^3]_0^{\pi/4}
  =\dfrac{1}{3}\left(1-\dfrac{\sqrt{2}}{4}\right) \\
  %
  (\text{第三項})
  &=\int_0^{\pi/4}cd\theta=\dfrac{\sqrt{2}}{2}  \\
  %
  (\text{第四項})
  &=\int_0^{\pi/4}c\theta d\theta=[s\theta+c]_0^{\pi/4}
  =\dfrac{\sqrt{2}\pi}{8}+\dfrac{\sqrt{2}}{2}-1
\end{align}
であるから\cref{eq:5}に代入して，
\begin{align}
  v
  &=\dfrac{1}{2}\left[\dfrac{5\sqrt{2}}{12}-\dfrac{1}{3}\left(1-\dfrac{\sqrt{2}}{4}\right)-\dfrac{\sqrt{2}\pi}{8}+\dfrac{\sqrt{2}\pi}{8}+\dfrac{\sqrt{2}}{2}-1 \right]    \\
  &=\dfrac{1}{2}\left(\sqrt{2}-\dfrac{4}{3}\right)
\end{align}
である．これを\cref{eq:3}に代入して
\begin{align}
  V=8\left(\sqrt{2}-\dfrac{4}{3}\right)r^3
\end{align}
である．$\cdots$(答)
\newpage
\end{document}
